\chapter*{\center \Large Abstract}

Large language models can write Verilog, but raw generation is an open loop: the
model never executes its own code. This report presents \textbf{SiliconScribe}, a
browser-based EDA playground that closes that loop with a real simulator. A natural
language prompt is parsed into a structured specification, which the user may review
and edit; RTL and a self-checking testbench are generated \emph{independently} from
that specification; Icarus Verilog 12.0 compiles and runs them; and failures are
sent back to the model, which must return a structured diagnosis and a corrected
design, for up to five iterations. The simulator decides PASS or FAIL, never the
model. Around that loop the system measures how much a PASS is worth: testbench
strength by mutation testing, toggle coverage measured from the simulator's own
waveform, a second-opinion testbench when a repair rewrote the testbench, and an
optional separate model for the testbench. After a PASS, Yosys maps the design to the
SkyWater 130\,nm cell library and OpenSTA reports slack, maximum frequency and area as
a \emph{pre-layout estimate} (one corner, no wire delay, not sign-off). The model is
chosen per run from the OpenCode app, Cloudflare Workers AI or NVIDIA NIM, with keys
read only from the server environment.

In a live evaluation with a free model (DeepSeek V4.1 Flash), \textbf{14 of 15}
prompts completed with no offline fallback: \textbf{8 passed on the first attempt,
6 after self-repair and none failed}. Their testbenches caught \textbf{102 of 111}
injected bugs (92\%), and independent hand-written reference testbenches found all 14
designs correct, so there was no false PASS. The fifteenth prompt timed out three
times and is reported as no result. The tool itself is checked by \textbf{453 backend
tests} and \textbf{222 frontend tests}, a mutation sweep of the built-in designs
(27 of 27 bugs caught), a gate-level simulation check of the mapped netlist on 19
designs, and an automated accessibility audit with zero violations. The mapped netlist
is also compared with the RTL by a formal equivalence check (Yosys SAT): 16 of 18
designs were proven equivalent and the two FSMs were reported as not proven. The architecture, the
engineering failure modes found during validation and the limitations that remain
are documented, together with a concrete scaling path (Section~\ref{sec:scale}).

\vspace{0.5cm}
\noindent
\textbf{Keywords:} Large Language Models, RTL Generation, Verilog, Functional Verification, Self-Checking Testbench, Icarus Verilog, Mutation Testing, Static Timing Analysis, Formal Equivalence Checking, Electronic Design Automation.
